\startsectionlevel[title={{\bf 90. Japanese Punctuation: How it works.}},reference={japanese-punctuation-how-it-works.}]

\goto{{\bf Japanese Punctuation: How it REALLY works. Lesson 90}}[url(https://www.youtube.com/watch?v=LDQlGnqElTU&list=PLg9uYxuZf8x_A-vcqqyOFZu06WlhnypWj&index=92&ab_channel=OrganicJapanesewithCureDolly)]

こんにちは。

Today we're going to talk about Japanese punctuation.

{\bf Now, the first thing we need to know about about Japanese punctuation is that most of it is a relatively late arrival in Japanese.} {\bf Most of it came in during the Meiji era, about a hundred and fifty years ago.}

{\externalfigure[../media/image156.webp]}

{\bf This was when Japan was modernizing, and a lot of the reason it came in was to help in translating Western literature.}

\thinrule

Now, {\bf what this means is that in most cases Japanese punctuation doesn't have such fixed and structural meanings as the equivalent punctuation does in English and other European languages.} So, we need to know what the punctuation is doing and also what it isn't doing.

\startsectionlevel[title={The full stop / 。},reference={the-full-stop}]

So, the first mark we're going to look at is the Japanese full stop or まる / 。, {\bf which looks like a little circle at the foot of a letter.}

{\externalfigure[../media/image193.webp]}

And {\bf this is really the exception to the rule, because it really is structurally clear.} {\bf What it does is end a sentence.}

{\externalfigure[../media/image731.webp]}

{\bf And this is very important for us in understanding how complex Japanese sentences are structured.}

Why is that? Well, {\bf Japanese is a very modification-heavy language.} {\bf That's to say, a lot of what other languages do by other strategies, Japanese does by modification}, that's to say, {\bf using clauses to modify nouns or other elements of a sentence.}

{\bf Most of the heavy lifting in Japanese structure that's not done by logical particles is done by this modification structure.}

{\bf So whole logical clauses can be used not as logical clauses but in order to modify a noun or some other element.}

So, let's take a look at how this works.

{\bf 市場で買って川に落としちゃった}お菓子。

And this means The candy {\bf that I bought at the market and done dropped in the river}.

I guess, done in the translation is used to point out the more literal translation of ちゃった.

{\externalfigure[../media/image589.webp]}

Now, {\bf the problem with this when we're reading a complex sentence is that we can get confused about whether something is a logical clause or not.}

So, {\bf are we saying here that I bought something at the market? No, we're not saying that.} {\bf Are we saying that I dropped something in the river?\crlf
No, that's not what we're saying in this sentence.}

\thinrule

{\bf Both of those things are just modifying お菓子.} {\bf They're telling us what kind of an お菓子 it was:} {\bf the one I bought at the market and done dropped in the river.}

{\bf So all we have here is a single noun, お菓子, that's been heavily modified by logical clauses that aren't working as complete logical clauses.}

\thinrule

{\bf We can then add a logical particle to that お菓子 and make it into a whole logical clause.}

So, we might say, 市場で買って川に落としちゃったお菓子{\bf が}魚さんに食べられた。

The candy{\em ({\bf =subject})} (I) bought at the market and done dropped in the river got eaten by a fish.

{\externalfigure[../media/image505.webp]}

If I understand correctly, because we marked お菓子 with が subject, we can now build another clause out of it, hence we get another logical clause of お菓子が魚さんに食べられた.

{\em While お菓子 is also modified by modifier clauses before it, specifically 市場で買って川に落としちゃった, they serve here as simply a modifier that describes お菓子 rather than being complete (standalone?) logical clauses. Although just my guess from how I understand it.}

{\bf But how do we ensure when we're reading a sentence like this whether we're looking at a logical clause or whether we're looking at a modifier?}

\stopsectionlevel

\startsectionlevel[title={How to know if it is a logical clause or a modifier?},reference={how-to-know-if-it-is-a-logical-clause-or-a-modifier}]

Now, {\bf the important thing to understand here is that in any kind of Japanese writing}\crlf
(except perhaps on Twitter or something)\crlf
{\bf a logical clause has to end in one of two ways}, {\bf either with a まる / 。, which tells us that it's the end of the complete sentence}

(it may have a couple of sentence ender particles after it, but it's a very strict rule of Japanese that {\bf what ends any logical clause is the B-engine, aside from any sentence-ender particles})

\thinrule

{\bf or it has to end in a clause connector that completes the clause and leads into the next clause.} {\bf This may be the て-form, it may be a word like から or けれど,} but armed with that information we can then see what's going on.

So, let's look at 市場で買って川に落としちゃった.

{\bf That's a complete pair of logical clauses.}

My guess here is that it is theoretically composed of two logical clauses if, like Dolly says below, we presume there is some hidden subject and a hidden direct object if we took both verbs as marking their individual logical clauses - 市場で(私がお菓子を)買った and (私がお菓子を)川に落としちゃった that are then connected using the て-form. Here, お菓子 should be a direct object hence why を, since 落とす and 買う are other-move verbs & require a Subject acting on a Direct Object. Obviously, topic marker can be there too, as 私は in the beginning.

{\em But instead, these clauses are combined and since they are followed by a noun serve instead as a modifier for that noun, so therefore they are not logical clauses anymore and do not contain the hidden elements, where instead they are one single big modifier for a noun that is following them, in this case お菓子. It is a shame I cannot ask Dolly about this, as I feel it would help a ton\crlf
However, I am aware my guess might very likely be wrong, so please take it with a lot of salt.}

{\bf It's not telling us what it was that we bought in the market and dropped in the river, but it could perhaps be presuming that.}

{\bf But we know that it's not in fact a logical clause because it doesn't end in a clause connector and it doesn't end with a まる / 。.} {\bf It goes straight into a noun.}

\thinrule

{\bf And this is how we can tell the difference between a modifier and a logical clause.} {\bf And we know when the entire sequence, the entire sentence, is over because it's going to have that まる / 。.}

And I've done a video about this method of analyzing Japanese sentences and I'll link that so that you can follow it up afterwards\goto{{[}34{]}}[url(./34-understand-any-sentence.md)].

{\bf But the まる / 。 here is vital.}

\stopsectionlevel

\startsectionlevel[title={The comma / 、},reference={the-comma}]

Now, the next one we're going to look at is the Japanese comma, which looks like a little diagonal line at the foot of a letter.

{\externalfigure[../media/image1054.webp]}

{\bf And that's kind of the opposite of a まる / 。, because it really is not a logical element in the sentence.} {\bf It was adopted into Japanese, but unlike the English comma or other European commas, it has no logical rules.}

{\bf Ideas like sectioning off a subordinate clause with commas at each end, that doesn't exist in Japanese.} {\bf You just put in a comma wherever you want to indicate a pause. And that's all it does.}

In Japanese schools, students are discouraged from using too many commas, and the reason for this is not incorrect comma usage, {\bf because there is no such thing as incorrect comma usage in Japanese.}

{\bf The reason is that commas are not a structural part of Japanese, so students are discouraged from using them as a crutch in conveying their meaning.} {\bf If you can't convey your meaning without commas, then you're not writing good Japanese.}

\stopsectionlevel

\startsectionlevel[title={The question mark / ？(はてな)},reference={the-question-mark-はてな}]

Now, the next one we'll look at is the question mark or はてな.

{\externalfigure[../media/image1109.webp]}

Now, {\bf in English the question mark is subject to a rule.} {\bf If you have a question, you have to end it with a question mark.}

{\bf And also in English, questions are structurally different from statements.}

So if we say, The coffee is hot, that's a statement, but if we say, Is the coffee hot? that's a question.

\thinrule

{\bf In Japanese we don't have this differentiation.} {\bf Statements and questions are structured exactly the same.}

Now, {\bf in formal Japanese} we have {\bf the question-marking particle か}, {\bf which tells us that it's a question.}

But {\bf in informal Japanese we can use か but mostly we don't.} {\bf Sometimes we use the question-marker の, but that's ambiguous because の can also be a statement marker.}

\thinrule

{\bf The only way you can really tell a question from a statement in Japanese is the rising intonation, which of course you can't hear in text.}

{\bf And so the question mark has become a very useful tool for indicating that rising intonation, for telling us that this is a question, not a statement.}

{\externalfigure[../media/image579.webp]}

{\bf In English, you have to use the question mark at the end of a question if you're writing proper English. In Japanese, there's no such rule.}

{\bf If the か marker is there, you don't use the question mark, but if you want to disambiguate the fact that something is a question rather than a statement, you just pop in a question mark if you want to.}

\stopsectionlevel

\startsectionlevel[title={The quotation marks / 「　」},reference={the-quotation-marks}]

The next thing we'll look at is quotation marks.

{\bf These look like little square brackets at the ends of a statement, and they work exactly the same as English quotation marks.}

{\externalfigure[../media/image956.webp]}

{\bf They just tell us that something is a quotation: it's what somebody's saying.} {\bf We don't use them for what somebody's thinking, as we sometimes do in English.}

\thinrule

Now, sometimes you'll see {\bf double quotation marks} like this, 『　』 {\bf and what they do is mark a quotation that occurs inside another quotation.}

\stopsectionlevel

\startsectionlevel[title={The side-marks},reference={the-side-marks}]

Now, the last thing I want to talk about is something that really puzzles people quite a lot.

{\bf In Japanese, particularly Japanese books}, you'll sometimes see, {\bf especially in vertical text, a set of little marks that look a bit like Japanese commas, running along the left-hand side of a word or a phrase.}

{\externalfigure[../media/image1065.webp]}

What on earth is this doing?

Nobody seems to tell you.

{\bf What it is actually doing is either emphasizing that word or phrase or telling us that it's being used in a special sense.}

\thinrule

{\bf So what these little marks are really doing is something like putting something into italics in English.}

And I suspect they came into Japanese in the first place to render italicization when translating Western literature.

{\bf There's no way of really writing Japanese characters in italics.} {\bf There is no italic format for Japanese, so this is what's used instead.}

\thinrule

{\bf This emphasis and the indication that something's being used in a special sense can also be indicated by katakana, but you will see this from time to time in Japanese texts and that's what it means.}

{\externalfigure[../media/image559.webp]}

{\externalfigure[../media/image260.webp]}

{\externalfigure[../media/image437.webp]}

{\em Dolly's link is for Lesson 80.}

\stopsectionlevel

\stopsectionlevel
